\subsubsection{Averaged Gradient Episodic Memory (A-GEM)}\label{sec:agem}

\gls{agem} \cite{AGEM} is a computationally efficient extension of
\gls{gem} that replaces the per-task \gls{qp} with a single average gradient
constraint. As with \gls{gem}, \gls{agem} requires task identity during training
to maintain and sample from per-task episodic memory buffers, but imposes no
such requirement at inference time. It is therefore equally compatible with both
\gls{til} and \gls{cil} evaluation settings.

\paragraph{Average Gradient Constraint}

Rather than enforcing \eqref{eqn:gem_constraint} independently for each prior
task, \gls{agem} constructs a single reference gradient $g_\text{ref}$ by
averaging over a batch of samples drawn uniformly from the episodic memories of
all previously observed tasks:
%
\begin{equation}\label{eqn:agem_ref}
    g_\text{ref} = \nabla_\theta \mathcal{L}(\theta; \tilde{\mathcal{M}})
\end{equation}
%
where $\tilde{\mathcal{M}}$ denotes a randomly sampled subset drawn from the
union of all prior task memory buffers. The proposed gradient $g$ is accepted
without modification if:
%
\begin{equation}\label{eqn:agem_constraint}
    \langle g, g_\text{ref} \rangle \geq 0
\end{equation}
%
If this constraint is violated, $g$ is decomposed into its components parallel
and orthogonal to $g_\text{ref}$ (Section~\ref{sect_bg_la}), and the parallel
component is discarded:
%
\begin{equation}\label{eqn:agem_projection}
    \tilde{g} = g - \frac{\langle g, g_\text{ref} \rangle}
    {\langle g_\text{ref}, g_\text{ref} \rangle} \, g_\text{ref}
\end{equation}
%
This is precisely the orthogonal complement projection $P^\perp g$ with respect
to the subspace spanned by $g_\text{ref}$, retaining only the component of $g$
that is orthogonal to $g_\text{ref}$.
%
This closed-form projection replaces the \gls{qp} of \gls{gem}, reducing the
per-step computational overhead from $\mathcal{O}(t)$ constraint evaluations to
a single inner product check regardless of the number of tasks observed.

\paragraph{Trade-off with GEM}

The efficiency gain of \gls{agem} comes at the cost of a weaker guarantee.
\gls{gem} ensures that the loss on \emph{every} prior task does not increase,
whilst \gls{agem} only constrains the loss averaged across prior tasks. In
practice, this means individual task performance may regress under \gls{agem}
even when the average constraint is satisfied. Nevertheless, \gls{agem} has
been shown to achieve competitive average accuracy to \gls{gem} at substantially
lower computational cost \cite{AGEM}, making it more tractable for longer
task sequences.